\thispagestyle{empty}
\chapter*{Abstract}
\addcontentsline{toc}{chapter}{Abstract}

\noindent\large
This report presents \textbf{AnsibleGuard}, a static-analysis system designed
and implemented from scratch for Ansible Infrastructure as Code during a
summer 2026 internship at the National School of Computer Science (ENSI),
University of Manouba. The system combines tolerant YAML/Jinja parsing,
artifact classification, rule-specific analysis, repository-level symbol and
execution graphs, and specialist evidence from ansible-lint and GitLeaks. It
supports full repositories, local changes, commits, and pull requests, and
exports a versioned, schema-validated output contract for reproducible analysis.

The current v0.0.7 study freezes 207 repositories. Its production run has 206
analyzable outputs, 35 partial ansible-lint runs, and one persistent repository
timeout; the 206 outputs contain 376,271 canonical occurrences. A paired
role-contract correction suppresses 247 unsupported variable findings on the
same analyzable population. A separate, corrected 207-repository threshold
benchmark extracts 2,986,459 unique observations for 24 metrics and prepares
16,971 review rows. The earlier 149-repository v0.0.6 study remains a
controlled historical comparison. Internal rules and targeted profiles were
designed and implemented in this project; ansible-lint and GitLeaks supply
clearly attributed specialist evidence. One agent-labelled variable review
supports calibration, but independent human labels and agreement are pending.
The report therefore distinguishes measured workload and semantic improvements
from unestablished population-wide precision and recall.

\medskip
\noindent\textbf{Keywords:} Infrastructure as Code, Ansible, Static Analysis,
Semantic Analysis, Code Review, Reproducible Evaluation.

\vspace{0.5cm}\noindent\rule{\textwidth}{0.4pt}\vspace{0.3cm}
\chapter*{Résumé}

\noindent\large
Ce rapport présente \textbf{AnsibleGuard}, un système d'analyse statique conçu
et développé depuis zéro pour le code d'infrastructure Ansible. Il associe
l'analyse tolérante de YAML/Jinja, la classification des artefacts, un graphe
d'exécution et de symboles, des règles internes, ainsi que les preuves fournies
par ansible-lint et GitLeaks. L'étude actuelle v0.0.7 porte sur 207 dépôts
figés. L'analyse de production a fourni 206 résultats exploitables et
376 271 occurrences canoniques; un dépôt a dépassé le délai et 35 exécutions
ansible-lint sont partielles. Une correction contrôlée des contrats d'arguments
de rôle supprime 247 signalements de variables injustifiés. Un banc d'essai
distinct couvre les 207 dépôts, 2 986 459 observations uniques et 24 métriques
numériques. Les règles internes et les profils ciblés ont été créés dans ce
projet. L'étiquetage humain indépendant reste nécessaire pour estimer la
précision et le rappel.

\medskip
\noindent\textbf{Mots-clés :} Infrastructure as Code, Ansible, Analyse
statique, Analyse sémantique, Revue de code, Évaluation reproductible.
